\begin{frame}[fragile]{Busca linear: localizar a primeira ocorrência}
\begin{lstlisting}
int alvo = 30;
int encontrado = -1;
for (int i = 0; i < 5; i++) {
    if (leituras[i] == alvo) {
        encontrado = i;
        break;
    }
}
if (encontrado >= 0) {
    printf("Encontrado no indice %d\n", encontrado);
} else {
    printf("Valor ausente\n");
}
\end{lstlisting}
\small O valor $-1$ significa ``não encontrado'' e nunca deve ser usado como índice. A busca visita no máximo cinco elementos neste exemplo.
\end{frame}
